\begin{proposicion}[Uniform convergence]\label{prop:conv_unif} As $x\downarrow0$, $r_{R,x},r_{C,x},u_{D,x},u_{A,x}$ converge uniformly
to positive continuous representatives of $r_{R,0},r_{C,0},u_{D,0},u_{A,0}$, which satisfy the bounds in Proposition~\ref{prop:cotas}
and equations~\eqref{eq:filas}--\eqref{eq:diagonales} with $x=0$ and $b=b_0$. \end{proposicion}

\begin{proof} Denote the integrals in~\eqref{eq:filas}--\eqref{eq:diagonales} by $\Phi_f$, so that $r_R=1/(x+b\Phi_R)$,
$r_C=1/(x+b\Phi_C)$, $u_D=1/(1+b\Phi_D)$, and $u_A=1/(1+b\Phi_A)$, and let $\Phi_{f,0}$ denote the same integral formed from the
limits. In each fiber integral, every factor from a different family is evaluated at an affine coordinate whose slope has absolute
value one or two. No slope is zero, since the four families have distinct directions. We bound the difference of the products by
summing the differences of their factors, which are bounded by $M$. A change of variable then gives
\begin{equation}\label{eq:fibras_unif}
 \sup_{t\in I_f}|\Phi_f(t)-\Phi_{f,0}(t)|\le C\sum_{g\ne f}\|y_{g,x}-y_{g,0}\|_{L^1(I_g)},
\end{equation} where $y_g$ is the normalized factor of the family $g$ ($r_R$, $r_C$, $u_D$, or $u_A$) and $C$ depends only on $M$. This
estimate holds even before we know that the limits are continuous. The limiting integrals $\Phi_{f,0}$ are continuous because
translations are continuous in $L^1$ and the fiber endpoints move continuously, as in Lemma~\ref{lem:vacancias}. By
Proposition~\ref{prop:conv_L2}, $\Phi_f\to\Phi_{f,0}$ uniformly. The denominators are bounded away from zero, since
$b\Phi_R,b\Phi_C\ge\kappa^2/(2(1+b_+))$ and $1+b\Phi_D,1+b\Phi_A\ge1$. The rational formulas now give uniform convergence to
$1/(b_0\Phi_{R,0})$, and similarly for the other factors. These limits are continuous, positive, and equal almost everywhere to the
limits in $L^2$. \end{proof}

To apply the implicit function theorem, we use the Banach space $X=C(I_R)\oplus C(I_C)\oplus C(I_D)\oplus C(I_A)$, with unknowns
$z=(z_R,z_C,z_D,z_A)=(\log r_R,\log r_C,\log u_D,\log u_A)\in X$ and $\gamma=\log b$. For $(x,z,\gamma)\in\mathbb R\times
X\times\mathbb R$, set $\nu=b\,r_Rr_Cu_Du_A$ and \[
 E_R=xr_R+T_R\nu-1,\quad E_C=xr_C+T_C\nu-1,\quad E_D=u_D+T_D\nu-1,\quad E_A=u_A+T_A\nu-1,
\] \[
 n_R=\int r_R-1,\qquad n_C=\int r_C-1 .
\] These are real-analytic maps between Banach spaces, formed from exponentials, products, bounded linear projections, and integrals.
Since $\int T_R\nu=\int T_C\nu$, we have the exact redundancy $\int E_R-\int E_C=x(n_R-n_C)$. Let $Z\subset X$ denote the closed
subspace $\{\int e_R=\int e_C\}$ of codimension one, and let $P(e)=(e_R,\,e_C-\int e_C+\int e_R,\,e_D,\,e_A)$, which maps $X$ onto $Z$
and acts as the identity on $Z$. Consider \begin{equation}\label{eq:sistema}
 \mathscr E(x,z,\gamma)=(PE,n_R,n_C)=0\qquad\text{in }Z\oplus\mathbb R^2 .
\end{equation} If $PE=0$ and $n_R=n_C=0$, the redundancy gives $\int E_C=\int E_R=0$ and hence $E=0$, also for $x\ne0$.
Thus~\eqref{eq:sistema} is equivalent to $E=0$ together with the two normalizations. By Lemma~\ref{lem:vacancias} and
Proposition~\ref{prop:conv_unif}, both the maximizers ($0<x\le\frac12$) and the limit ($x=0$) satisfy this system.

\begin{teorema}[Analytic branch at the full board]\label{thm:rama} There exist $\varepsilon>0$ and a real-analytic map
$x\mapsto(z(x),\gamma(x))$ from $(-\varepsilon,\varepsilon)$ to $X\times\mathbb R$ that solves~\eqref{eq:sistema}, with \[
 (z(0),\gamma(0))=(\log r_{R,0},\log r_{C,0},\log u_{D,0},\log u_{A,0},\log b_0),
\] and that, for $0<x<\varepsilon$, coincides with the optimal profile: \[
 (z(x),\gamma(x))=(\log r_{R,x},\log r_{C,x},\log u_{D,x},\log u_{A,x},\log b_x).
\] \end{teorema}

\begin{proof} For a variation $(h,\beta)=(h_R,h_C,h_D,h_A,\beta)\in X\oplus\mathbb R$, set $w=\beta+\sum_fh_f\circ\ell_f$, so that
$\dot\nu=\nu_1w$ at $x=0$. The part of the derivative of $E_f$ acting on $h_f$ is $h_f$. For rows and columns this follows from
$T_R\nu_1=T_C\nu_1=1$, and for a diagonal it follows from $(u_{D,0}+m_D[\nu_1])h_D=h_D$. The operators coupling distinct families are
compact by Arzelà--Ascoli. When parametrized by the coordinate of the other family, their kernels are continuous on the domain allowed
by the fiber (extending them by zero may introduce jumps at the boundary of that domain, which does not affect the argument). For two
nearby fiber parameters, we use uniform continuity of the kernel on the common domain, and the symmetric difference of the intervals of
integration has length $O(|t-t'|)$. For $D\leftrightarrow A$, the change of variable has Jacobian $\frac12$ and endpoints
$[|d|,2-|d|]$. The term in $\beta$ has rank one. Thus the derivative of $E$ is a Fredholm operator from $X\oplus\mathbb R$ to $X$ of
index one. At $x=0$ its image lies in $Z$ because the derivative of the redundancy vanishes. This raises the index to two when the
operator takes values in $Z$. The two normalization equations bring the index back to zero.

If $(h,\beta)$ lies in the kernel of the derivative of~\eqref{eq:sistema} at $x=0$, then \begin{gather*}
 T_R(\nu_1w)=T_C(\nu_1w)=0,\qquad T_D(\nu_1w)=-u_{D,0}h_D,\qquad T_A(\nu_1w)=-u_{A,0}h_A,\\
 \int r_{R,0}h_R=\int r_{C,0}h_C=0 .
\end{gather*} Multiplying by the components of $w$, integrating, and using the row equation to eliminate $\beta\int\nu_1w=\beta\int
T_R(\nu_1w)=0$ gives \begin{equation}\label{eq:energia}
 \int_{[0,1]^2}\nu_1w^2+\int_{I_D}u_{D,0}h_D^2+\int_{I_A}u_{A,0}h_A^2=0 .
\end{equation} By positivity, $h_D=h_A=0$ and $\beta+h_R(r)+h_C(c)=0$ on the square. Hence $h_R$ and $h_C$ are constant, the
normalizations force them to vanish, and then $\beta=0$. The kernel is trivial, and since the index is zero, the derivative is an
isomorphism with bounded inverse. The analytic implicit function theorem~\cite{Zeidler} gives a neighborhood $U$ of
$(0,z(0),\gamma(0))$ and an analytic branch whose points are precisely the solutions of~\eqref{eq:sistema} in $U$. By
Proposition~\ref{prop:conv_unif} and the lower bounds, the logarithms of the optimal profiles converge uniformly as $x\downarrow0$. For
small $x>0$ they lie in $U$ and coincide with the branch. \end{proof}

Propositions~\ref{prop:cotas}, \ref{prop:conv_L2}, and~\ref{prop:conv_unif} and Theorem~\ref{thm:rama} prove Theorem~\ref{thm:lleno}.
